% --- 1. ОБЩИЕ НАСТРОЙКИ ТЕКСТА ---
\onehalfspacing
\setlength{\parindent}{1.25cm} 

% Улучшенное выравнивание по ширине и переносы
\sloppy 
\hyphenpenalty=1000 
\tolerance=2000

% --- 2. ПРИНУДИТЕЛЬНОЕ "СОДЕРЖАНИЕ" (Название) ---
\usepackage{etoolbox}
\addto\captionsrussian{%
  \renewcommand{\contentsname}{СОДЕРЖАНИЕ}%
}

% --- 3. ЦЕНТРИРОВАНИЕ ЗАГОЛОВКА СОДЕРЖАНИЯ (И только его!) ---
% Мы переопределяем команду вывода заголовка списка (оглавления)
\renewcommand*{\deftocheading}[2]{%
  \begin{center}
    \bfseries\normalsize #1
  \end{center}
  \vspace{1em}
}

% --- 4. ШРИФТЫ ЗАГОЛОВКОВ И СОДЕРЖАНИЯ (Везде 14pt) ---
\setkomafont{disposition}{\rmfamily\bfseries\normalsize}
\setkomafont{chapter}{\rmfamily\bfseries\normalsize}
\setkomafont{section}{\rmfamily\bfseries\normalsize}
\setkomafont{subsection}{\rmfamily\bfseries\normalsize}

% Пункты в оглавлении — обычные, не жирные
\setkomafont{chapterentry}{\normalfont\normalsize}
\setkomafont{chapterentrypagenumber}{\normalfont\normalsize}
\KOMAoptions{toc=chapterentrydotfill, numbers=noendperiod}

% --- 5. МАГИЯ ОТСТУПОВ ЗАГОЛОВКОВ (1.25 см) ---
\makeatletter
% Для разделов (section, subsection)
\renewcommand{\sectionlinesformat}[4]{%
    \hspace{1.25cm}#3#4%
}

% Для глав (chapter)
\renewcommand{\chapterlinesformat}[3]{%
    \hspace{1.25cm}#2#3%
}
\makeatother

% Интервалы
\RedeclareSectionCommand[beforeskip=12pt, afterskip=6pt]{chapter}
\RedeclareSectionCommand[beforeskip=12pt, afterskip=6pt]{section}
\RedeclareSectionCommand[beforeskip=12pt, afterskip=6pt]{subsection}

% --- 6. СПИСКИ (enumitem) ---
\setlist{
	leftmargin=0pt,
	nosep,
	wide=\parindent,
	labelsep=0.5em,
	listparindent=\parindent
}
\setlist[itemize,1]{label=--} 
\setlist[enumerate,1]{label=\arabic*.}

% --- 7. ПОДПИСИ И ТАБЛИЦЫ ---
\captionsetup{justification=centering, singlelinecheck=false}
\DeclareCaptionLabelFormat{simple}{#1~#2}
\captionsetup{labelformat=simple}
\DeclareCaptionLabelSeparator{emdash}{ --- }
\captionsetup[table]{
	labelsep=emdash,
	justification=raggedright,
	singlelinecheck=false,
	position=top,
	skip=6pt
}

% --- 8. ЛИСТИНГИ ---
\lstset{
    basicstyle=\ttfamily\footnotesize,
    breaklines=true,
    texcl=true,
    frame=none
}